\begin{table}[H]
\small
\setlength{\tabcolsep}{5pt}
\renewcommand{\arraystretch}{0.95}

\caption{\label{tab:country-heterogeneity-cumulative-gap}Individual P20--P80 inflation ranges and Q1--Q5 inflation inequality in the euro area, 2021--2023}
\centering
\begin{tabular}[t]{llllllll}
\toprule
\multicolumn{2}{c}{} & \multicolumn{6}{c}{Inflation inequality}\\
\cmidrule(lr){3-8}
\multicolumn{2}{c}{} & \multicolumn{3}{c}{Individual P20--P80} & \multicolumn{3}{c}{Income Q1--Q5}\\
\cmidrule(lr){3-5} \cmidrule(lr){6-8}
Country & \shortstack{Weight\\(2022, \%)} & 2021 & 2022 & 2023 & 2021 & 2022 & 2023\\
\midrule
Germany & 28.1 & 1.4 & 4.3 & 1.8 & -0.4 & 1.4 & 0.5\\
France & 20.4 & 1.5 & 2.6 & 2.2 & 0.4 & 0.2 & 0.2\\
Italy & 16.3 & -- & -- & -- & 0.1 & 1.6 & 0.5\\
Spain & 11.0 & 2.5 & 4.0 & 3.6 & 0.8 & 0.9 & -0.7\\
Netherlands & 5.4 & 2.0 & 6.8 & 3.7 & 0.0 & 3.1 & -1.2\\
\addlinespace
Belgium & 4.0 & 3.3 & 8.0 & 7.3 & 1.3 & 2.3 & -2.7\\
Austria & 3.3 & 1.2 & 3.8 & 2.2 & -0.5 & 0.0 & 0.5\\
Portugal & 2.2 & -- & -- & -- & 0.8 & 0.2 & -0.9\\
Greece & 2.2 & 2.1 & 4.7 & 2.7 & 0.3 & 1.2 & -0.6\\
Finland & 1.9 & 1.9 & 3.7 & 2.1 & -0.2 & -0.9 & 0.4\\
\addlinespace
Ireland & 1.5 & 2.0 & 5.0 & 3.1 & 0.5 & 2.9 & 1.1\\
Slovakia & 0.8 & 1.9 & 2.6 & 2.6 & -0.9 & 0.1 & 0.5\\
Croatia & 0.7 & 1.2 & 2.7 & 2.5 & 0.0 & 0.2 & -0.3\\
Lithuania & 0.6 & 2.2 & 10.1 & 4.4 & 0.2 & 3.3 & 0.4\\
Slovenia & 0.4 & 1.4 & 2.4 & 2.3 & 0.0 & 0.4 & 1.5\\
\addlinespace
Luxembourg & 0.3 & 2.6 & 4.9 & 2.8 & 0.1 & 0.2 & -0.1\\
Latvia & 0.3 & 1.5 & 7.6 & 3.8 & -0.2 & 4.5 & 2.5\\
Estonia & 0.2 & 3.0 & 10.7 & 3.4 & 0.6 & 4.8 & 1.3\\
Cyprus & 0.2 & 1.5 & 3.0 & 1.7 & 0.0 & 0.8 & 0.2\\
Malta & 0.1 & 1.9 & 2.2 & 2.2 & -0.2 & 0.3 & 0.6\\
\addlinespace
\midrule
\textbf{Euro area} & \textbf{100.0} & \textbf{1.7} & \textbf{4.2} & \textbf{2.6} & \textbf{0.1} & \textbf{1.2} & \textbf{0.0}\\
\bottomrule
\end{tabular}
\vspace{-0.4em}
\begin{flushleft}
\footnotesize{Notes. The first block is the P80 minus P20 range of annual household inflation, in percentage points, using national HBS survey weights. Italy and Portugal are unavailable at the required household-product detail and are shown as dashes; the euro-area row averages national endpoints over the other 18 countries using annual HICP country weights. The second block is the Q1 minus Q5 gap for all 20 countries. Annual rates compare annual-average indices with the preceding year. The displayed country weights are 2022 EA20 HICP shares. Sources: Eurostat, HBS 2020 microdata, authors' calculations.}
\end{flushleft}
\end{table}
